\scanpage{212}
\mspage{161}

\uline{Corollaire 2} Soit $Q$ un carré cartésien (p. 158), avec $i'$ (donc \inserted{aussi} $i$) une imm. ouverte, et $q$ \inserted{(donc aussi $p$)} un \uline{isom.} \uline{local}. \deleted{Supposons que \unsure{$T$}} \inserted{Soit $T=X\setminus U$, $T'=X'\setminus U'$} (donc $T\simeq X\times_{X'}T'$), et supposons $T'\longrightarrow T$ iso. Alors $Q$ est cocartésien.

C'est un cas particulier du cor. 1, mais sans référence explicite à un localiseur $W$. On peut prendre $W=W_\infty$, ou $W=W_0$, pour se ramener au corollaire 1. Mais il est plus raisonnable de dire que dans le carré cartésien $(*)$, l'hypothèse que $q$ soit un iso. local implique que $q_{/t}$ \unsure{est} un isomorphisme, donc $p_{/t}$ l'est aussi. Il en résulte bien que les $v'\backslash(U_{/t})$ sont $0$-connexes, \deleted{\unsure{si ce}} \deleted{\unsure{puis}} \unsure{cette|ces} \unsure{catégorie|catégories} \unsure{étant} isom. à $\unsure{v'(U'_{/t'})|v'\backslash(U'_{/t'})}$, donc a un objet initial.

\deleted{Consid} Soit maintenant $X$ un objet de Cat, et considérons un diagramme cocartésien dans $X^\wedge$ :

\diagram{\begin{tikzcd}[ampersand replacement=\&]
\& \unsure{G|G'}\arrow[r,hook,"i"]\arrow[d,"p"'] \& \unsure{F|F'}\arrow[d,"q"]\\
(\unsure{*|**})\& G'\arrow[r,hook,"i'"] \& F',
\end{tikzcd}}
% archive sheet number: 212
